（五）怎样解释词语

有人怕解词，说“词语那么多，在不同语言环境里词义变化又那么大，我们怎么能全记下来呢？”

是的，要学好语言，非下苦功不可。但是，这并不是说必须去背词典。其实，死搬词典的解释，有时还讲不通具体文章中的某个词呢。比如，朱自清先生的著名散文《荷塘月色》中有句话：“叶子底下是脉脉的流水。”《现代汉语词典》解释“脉脉”是“默默地用眼神或行动表达情意”的意思，如“脉脉含情”。然而这儿的“流水”不是人，无所谓“眼神或行动”。原来朱先生在此用了拟人的修辞手法，巧
\newpage
\noindent
妙地“形容水没有声音，好象深含感情的样子。”如果照抄词典的解释，是要碰壁的。

这儿有条解词的原则我们要记住，那就是“字不离词，词不离句，句不离段。”只要我们把课文中的词语记牢，再掌握一些关于词语的基本知识，做到日积月累，举一反三，就能逐渐地把知识化为能力，就能基本正确地理解一般文章中的词语。

具体说来，我们在解词时要注意些什么问题呢？

要弄清词素。现代汉语中的词大多数是双音节合成词，这些词的每个词素（具有最小意义的构词单位）都同词义有关。比如，“庐舍”的“庐”指简陋的房屋，“舍”也是房屋，因此这个词就有“房屋、田舍”的意思了。再如，“探访”这个词就是“探寻、访求”的意思。其他如“糟粕”、“遏制”、“清晰”、“精密”、“阻塞”、“智慧”等等都是由意义相同或者意义相近的词素组成的合成词，我们只要弄清了每个词素的含义，整个词义基本上都能迎刃而解。对于由意义相反或意义相对的词素组成的合成词，要解释它们的词义，也应从弄清词素着手。例如：“毁誉”的“毁”指“毁谤”，“誉”指“称赞”，整个词义正是“毁谤和称赞”的意思。其他如“兴亡”、“旦夕”、“曲直”、“利害”、“出纳”、“好恶”、“是非”、“动静”等等的词义，都是由它们词素的意义构成的，只不过在具体语言环境中有着特定的引申义或比喻义罢了。比如“旦夕”，指早晨和晚上，比喻短时间内；再如“是非”，指事理的正确和错误，引申为“口舌”，因说话而引起的误会或纠纷。至于辨析同义词的词义，我们更要弄清它们词素

\newpage
\noindent

的异同，从而作出准确的解释。例如：“屹立”、“耸立”、“挺立”、“矗立”、“壁立”，相同的词素是“立”，说明它们都有“站立”的意思；不同的词素是“屹”（yì，象山峰一样高耸而稳固，常用来比喻坚定不可动摇）、“矗”（高耸而直）、“挺”（硬而直，常含有坚强不屈的意味）、“耸”（高起）、“壁”（象墙壁一样陡），我们弄清了它们的不同词素的含义，就能对这些词作基本准确的解释和正确的运用。

根据这种辨析词义的方法，我们就可以把由同义或反义的词素构成的合成词的词义，大致正确地解释出来，例如：“丰满”是说丰腴、饱满，指胖得匀称好看；“恶劣”是说很坏；“荒诞”是说荒谬、怪诞，指极不真实、极不近情理；“好歹”就是好坏，引申有不管怎样的意思。我们也可以根据辨析词义的方法，把成语中的同义词素或反义词素的含义反推出来。例如“文过饰非”（“文”可由“饰”反推出“掩饰”的意思），怨天尤人（“尤”可由“怨”反推出“怨恨”的意思），求全责备（“责”可由“求”反推出“要求”的意思）；前倨后恭（“倨”可由“恭”反推出“傲慢”的意思），苦尽甘来（“甘”可由“苦”反推出“甜”的意思），一张一弛（“弛”可由“张”反推出“松懈”的意思），等等。另外，我们也可通过对同义词中的不同词素的辨析，区别一组词的词义。例如“安静、平静、宁静、幽静”，“精巧、精细、精美、精致”，“严格、严明、严正、严厉”，“改进、改善、改良”，“推辞、推托、推卸”，等等。

要弄清各词素间的结构方式。我们前面已经说过，
\newpage
\noindent
现代汉语中的词，按照结构特点，可以分为单纯词和合成词两类。单纯词是由一个词素构成的词，它是一个整体，不论这类词素包含几个字，都不能拆开用、拆开讲，比如“志”、“逍遥”、“踽踽”、“蹉跎”、“窈窕”、“踽踽”、“玻璃”、“梵哑铃”、“布尔什维克”等等。而合成词大多数是由两个词素构成的，它们的结构方式主要有：联合式（约略、贫瘠、反正、动静等等）、偏正式（火花、流俗、肆虐、偏信等等）、主谓式（雷鸣、蚕食、鲸吞、眼红等等）、动宾式（匿名、掉臂、揭幕、散心等等）、述补式（改善、打倒、压扁、提高等等）。我们在解词时，只要根据语言环境弄清了词的结构方式，就比较好办了。例如：

（1）这些无视流俗的新青年，唱着歌儿，掉臂而行。

这儿的“流俗”，是偏正式结构，指“流行的”（或者一般的）风俗习惯（含贬义），而不是联合式结构和动宾式结构，所以不能讲成“流言和习俗”或者“流掉旧习气”；“掉臂”是动宾式结构，指“甩着臂膀（走路）”，形容自由自在的样子（这儿的“掉”是“摆动”的意思），而不能解释为偏正式结构“掉下来的臂膀”，当然也不是“掉了臂膀”的意思。

（2）居里夫人在任何时候都意识到自己是社会的公仆。由于社会的严酷和不平等，她的心情总是抑郁的。

这儿的“意识”是主谓式结构，指“心里认识到（或者觉察到）”，属于动词，而不是指“人的头脑对于客观物质世界的反映”，不能用名词的意思来讲；“严酷”是联合式结构，有“严厉、冷酷”的意思，在这儿偏向“残酷、冷酷”的意思。

\newpage
\noindent

要弄清词性。现代汉语的某个词一般均有比较固定的词性，但是在不同的语言环境中，词性也是可以变化的。我们必须根据词不离句的原则分析词性，从而给以准确的解释。上面所说的“意识”一词就是一个明显的例子。再如廖承志《致蒋经国先生信》中说：“再次合作，老先生主其事”，“试为贵党计……”这儿的“主”和“计”都处于谓语位置，就不能用名词“主人”、“计谋”来解释，而应讲为动词“主持、主管”、“打算、筹划”才行。又如高考试卷中曾叫考生解释“世故深通”的“世故”一词，很明显，这儿是个名词，应作“处世的经验”讲，而与“此人很世故”中作动词讲的“圆滑，不得罪人”的词义有别。

根据这种辨析方法，我们就可以在弄清词语所处的不同语言环境的基础上，确定词性，解释词义。

要弄清词的基本义和转义。例如，“分”的基本义是“分开”，由这个基本义派生出来的转义有“分配”、“辨别”、“分支”、“分数”等等。为了表意的需要，人们在写文章时，不仅运用词的基本义，而且更广泛地运用词的转义（一般包括引申义和比喻义），这就要凭借具体的语言环境，联系句、段的文意进行辨析。例如：“重要的问题在善于学习。”“问题”的基本义是指“要求回答或解释的题目”，有时引申为“需要研究讨论并加以解决的矛盾、疑难”（如“工作中总会遇到问题”），而在这儿却指“关键、重要之点”，用的是由基本义派生出来的又一个引申义。再如，高考试卷中曾叫考生解释“四万万人坚决地沉着地接受炮火的洗礼了……我们以殉道者的精神负起我们应负的十字架！”这一句中的“洗礼”和“殉道者”都不再是基督教义
\newpage
\noindent

中的术语原义了，而是分别比喻“重大斗争的锻炼和考验”与“为维护正义、追求革命理想而牺牲生命的人”。这种词的比喻义用法，在成语中特别多，如“胸有成竹”、“满城风雨”、“养痈遗患”、“仁者见仁，智者见智”、“狐假虎威”、“死灰复燃”、“敝帚自珍”、“唇亡齿寒”等等，在解释这些词语时，一定要从词的基本义出发，领会它的整体意义，准确地讲出此词在本文中的比喻义，切不要望文生义，否则就讲不通整句的意思，有时甚至会闹出笑话来。

要灵活运用古代汉语中的词语知识。因为不少古代汉语的词汇，特别是成语，还活在现代汉语的词语中，所以在解词（特别是解释成语）时，就要融会贯通已经学到的古代汉语词汇的知识。例如，高考试卷中曾叫考生解释“自圆其说”。这里关键的一个字是“圆”，我们只要想到了古代汉语中的使动词用法，立即就可以把它解释成“使自己说话圆满，没有自相矛盾的地方”。再如让考生解释“敝屣尊荣”的“敝屣”。这是古代汉语意动词用法的一个例子，只要想到了这一点，词义就可以解释成“把……看成象破鞋子一样，比喻极端看不起”。其实，在高考试卷中，涉及古代汉语知识的词语很多，如“举国一致”、“付之一炬”、“殉道者”、“持之有故，言之有理”、“诚然”、“偶然”等，我们解释这些词，都应想到灵活运用古代汉语的词语知识才行。

当然，解词时还必须根据上下文意，确定词义的感情色彩，如“骄傲”一词究竟是“自以为了不起”（贬义词），还是“自豪”（褒义词）的意思？这要由上下文意来决定。

\newpage
\noindent
在一定的语言环境里，有些联合式的合成词只能取其中一个词素的含义，这叫偏义词。解释时要特别留心。如“无数革命先烈为了中华民族的生死存亡而浴血奋战”，这里的“生死存亡”只取“生存”的含义，而不取“死亡”的含义，是一目了然的。

我们基本掌握了一个词的含义，又如何进行表述呢？较常见的有如下几种方式：

（1）用同义词语解释。它的原则是：用好懂的同义词语解释不易懂的词语。例如：
驰骋：奔跑。（用常用词解释生僻词。）
对劲：合得来。（用普通话解释方言。）
愚不可及：愚蠢到了极点。（用现代语解释文言。）
水门汀：水泥。（用汉语解释外来语。）

（2）用反义词加否定词解释。它的原则是：这反义词非常明显而又为人们所熟悉。例如：
晦气：不吉利。（同义词语是“倒霉”。）
鬼鬼祟祟：不光明正大，不正派。
白话：对文言而说的一种书面语言。

（3）用分析字义综合词义的方式解释。这只适用于字义和词义之间的关系非常明显的合成词或成语。例如：
夙愿：夙，旧有的，一向有的；素有的愿望。
自诩：诩，夸耀；自我夸耀。
怒形于色：形，表露；色，脸色；愤怒显现在脸上。

（4）用说明事物的类属及其特点的方式解释。事物的特点包括其内容、性质、状态、功用等方面。例如：
七窍：人脸上的口、眼、鼻、耳共有七个孔，叫做“七

\newpage
\noindent

窍”。（从内容上说）
冒充：以假作真欺骗别人。（从性质上说）
立体：有长、宽、高的物体。（从状态上说）
勺（sháo）：舀东西的用具，略呈半球形，有柄。（从功用和状态上说）
耙：用来破土松土的农具。（从功用上说）

（5）运用形容、比喻的方式解释。形容一般要说出感情色彩，比喻应与要解释的对象相近。例如：
顶天立地：形容英雄气概。
铁面无私：形容公正严明，不讲情面。
棘手：比喻事情难办。
焦头烂额：比喻十分狼狈窘迫。
胸有成竹：比喻做事之前已经有通盘的考虑。

其它还有一些表述方式，如下定义等，不再赘述。
总之，为了准确地解释词语，必须把词语放到具体的句子中，联系上下文意，综合运用关于词语的知识和灵活运用解词的方式，才有可能做到基本理解和表达一个词的含义，才有可能提高阅读能力。
